\section{Conclusion}
HEAR-WORK AI provides a public-data-first, clinically grounded, and openly reproducible framework for studying tinnitus, hearing loss, and stress-related exposures in employed adults. KNHANES is the preferred development dataset, NHANES the preferred public external-validation dataset, and open medical AI models accelerate feature extraction and explanation only under schema constraints, a deterministic red-flag safety layer, and strict human verification. By prespecifying aims, honoring complex-survey design, reporting calibration and fairness, and releasing code, model cards, and a synthetic-data path, the framework improves reproducibility, robustness, and global usefulness while avoiding unsupported causal claims. Public data can establish associations and benchmark methods; the prospective clinical extension is required to study treatment timing and recovery---above all, to protect the narrow window in which sudden hearing loss is treatable.
